\documentclass[11pt]{scrartcl}
\usepackage[a4paper,margin=2.4cm]{geometry}
\usepackage{../common/rlstudy}

\title{\textbf{Week 20 --- Extended reading}\\\large Critic-free RL, RLVR, and process reward models}
\author{STUDY-RL}
\date{\today}

\begin{document}
\maketitle

\section*{How to use this list}
Five core papers, ranked roughly by how often the rest of the field cites
them. Each entry includes a one-paragraph summary, a \emph{why this matters
for week 20}, and an annotated key passage with two to three sentences of
commentary. Sourced from \texttt{research/llm\_rl\_sota.md} (sections 3--5)
and \texttt{research/safety\_rl\_sota.md} (section 5).

% ===========================================================================
\section{Shao et al. --- DeepSeekMath / GRPO}

\textbf{Citation.} Shao, Z., Wang, P., Zhu, Q., Xu, R., Song, J., Bi, X.,
\dots, Wu, Y. (2024). \emph{DeepSeekMath: Pushing the Limits of Mathematical
Reasoning in Open Language Models}. arXiv:2402.03300.

\textbf{Summary.} The DeepSeekMath paper introduces Group Relative Policy
Optimisation (GRPO). It samples $G$ completions per prompt, scores each with
a reward, and forms the advantage by group-mean subtraction and
group-standard-deviation division. The advantage is broadcast to every
token of the completion and plugged into a PPO-style clipped surrogate
with a KL-to-reference penalty --- crucially \emph{without} a value head.
On the 7B DeepSeekMath-Base, GRPO with rule rewards reaches a 51.7\%
top-1 on MATH, beating prior open-weights baselines and matching some
closed models.

\textbf{Why it matters for this week.} GRPO is now the dominant critic-free
algorithm for LLM RL --- every major 2025 reasoning-RL recipe (DeepSeek-R1,
Kimi K1.5, Open-Reasoner-Zero, DAPO) builds on its loss formulation. The
paper is also the cleanest derivation of the
group-baseline-as-control-variate idea applied to language models.

\begin{quotation}
\itshape ``The reason we choose to use such an objective function is that
the value function used in PPO is typically another model of comparable
size as the policy model, which brings substantial memory and computational
burden\ldots GRPO, instead, estimates the baseline from group scores,
significantly reducing training resources.''
\end{quotation}
This is the central thesis: the critic was paying twice --- once in
memory, once in fragility of the on-policy value fit --- and the group
mean is a strictly better baseline for end-of-sequence rewards because it
is unbiased by construction and uses no learned parameters.

% ===========================================================================
\section{Ahmadian et al. --- RLOO}

\textbf{Citation.} Ahmadian, A., Cremer, C., Gallé, M., Fadaee, M., Kreutzer,
J., Pietquin, O., \dots, Üstün, A. (2024). \emph{Back to Basics: Revisiting
REINFORCE Style Optimization for Learning from Human Feedback in LLMs}.
arXiv:2402.14740.

\textbf{Summary.} The paper argues that PPO is overkill for RLHF: token-level
advantages are unnecessary when the only reward signal is at the end of the
sequence, and the value head adds compute without adding signal. Their
proposal is REINFORCE-Leave-One-Out: sample $k$ completions per prompt,
use the average of the other $k - 1$ as the baseline for each, treat the
whole completion as a single action, and take a sequence-level policy
gradient step. On a 6.9B model, RLOO outperforms a heavily tuned PPO at
$3\times$ the throughput and $70\%$ less GPU RAM.

\textbf{Why it matters for this week.} RLOO is the conceptually cleanest
critic-free algorithm and the foundation on which GRPO sits. It is also
the reference point for variance-reduction analysis in lectures, because
its unbiasedness is exactly the LOO independence argument
(Theorem~\ref{thm:rloo-unbiased} in \texttt{notes.tex}).

\begin{quotation}
\itshape ``Treating PPO's clipping and advantage estimation as
\emph{required} features\ldots papers over the fact that, in the RLHF setting,
the variance of the simple REINFORCE estimator with a leave-one-out
baseline is small enough that the additional complexity of GAE,
value-network fitting, and per-token credit assignment is not warranted.''
\end{quotation}
The passage is also a methodological lesson: it is easy in deep RL to
inherit machinery from prior work without checking whether the
assumptions it was solving for still apply. The PPO clip and GAE are
brilliant in continuous-control settings with dense reward and long
horizons; for terminal-reward, single-action LLM RL they are largely
ornamental.

% ===========================================================================
\section{Liu et al. --- Dr.GRPO}

\textbf{Citation.} Liu, Z., Chen, C., Li, W., Qi, P., Pang, T., Du, C.,
Lee, W. S., Lin, M. (2025). \emph{Understanding R1-Zero-Like Training:
A Critical Perspective}. arXiv:2503.20783.

\textbf{Summary.} A careful audit of the original GRPO loss isolates two
biases: the $1/|y|$ token average penalises long correct chains-of-thought
relative to short accidental ones (the ``length bias''), and the
$1/(\sigma_r + \varepsilon)$ amplifies gradients on prompts whose group
has small intra-variance (the ``difficulty bias''). The Dr.~GRPO fix is to
sum rather than average over tokens, and to drop the standard-deviation
divisor. The patched loss closes the length-vs-quality decoupling that
many RLVR runs were silently exploiting.

\textbf{Why it matters for this week.} Dr.~GRPO is the example par excellence
of a debugging insight that a reader who only sees the algorithm via its
formula will miss. The two changes are tiny --- a denominator and a sum ---
but they shift the implicit regulariser of the loss in ways visible only
when you plot completion-length distributions over training.

\begin{quotation}
\itshape ``The standard GRPO objective implicitly defines a non-uniform
weighting over tokens and groups that bears no relation to the reward
landscape; correcting it is a one-line change and uniformly improves
both accuracy and chain-of-thought length stability.''
\end{quotation}
The deeper moral: an unbiased estimator at the \emph{population} level can
still be biased at the \emph{loss surface} level if the per-example
normaliser depends on quantities that correlate with the target. This
echoes the classical statistical lesson on weighting under heteroskedastic
noise.

% ===========================================================================
\section{DeepSeek-AI --- DeepSeek-R1}

\textbf{Citation.} DeepSeek-AI (2025). \emph{DeepSeek-R1: Incentivizing
Reasoning Capability in LLMs via Reinforcement Learning}. arXiv:2501.12948.

\textbf{Summary.} R1 is the production-scale instantiation of GRPO~+~RLVR.
The R1-Zero variant trains from a base (non-SFT) model purely on
rule-rewarded math, code, and logic tasks, with no preference data and no
demonstrations. Despite the cold start, the model develops a long
chain-of-thought style spontaneously --- the paper documents ``aha
moments'' where the model verifies its own derivations mid-stream. R1
proper adds a brief cold-start SFT, then multiple rounds of RL, beating
GPT-4-level baselines on AIME and MATH while remaining open-weights.

\textbf{Why it matters for this week.} R1 is the empirical demonstration
that the GRPO + verifiable-reward recipe scales. It also empirically
contradicts the conventional wisdom that one cannot do useful RL from a
base model: with a strong base, the rule reward alone provides enough
signal for emergent CoT.

\begin{quotation}
\itshape ``DeepSeek-R1-Zero showcases capabilities such as self-verification,
reflection, and generating long CoTs\ldots without any supervised data,
indicating that reasoning capability can emerge purely through RL,
provided the reward is verifiable.''
\end{quotation}
This is a strong claim. It rests on the base model already having latent
reasoning circuits that RL surfaces, rather than RL creating them from
scratch --- consistent with the weak-to-strong narrative in week~21.
The R1 verifier is a regex on the boxed final answer plus a code
sandbox: a strikingly minimal interface for such a strong result.

% ===========================================================================
\section{Lightman et al. --- PRM800K}

\textbf{Citation.} Lightman, H., Kosaraju, V., Burda, Y., Edwards, H., Baker,
B., Lee, T., Leike, J., Schulman, J., Sutskever, I., Cobbe, K. (2023).
\emph{Let's Verify Step by Step}. arXiv:2305.20050.

\textbf{Summary.} The paper releases PRM800K --- 800,000 human step-labels
on MATH problem solutions --- and trains a process reward model that
scores every intermediate step. At best-of-N inference, the PRM
re-ranker pushes a 6B base solver from $52\%$ to $78\%$ on MATH at
$N = 1860$, comfortably above the same model's ORM counterpart. The paper
also studies which step features the human labellers attended to, and
provides an early discussion of \emph{distributional shift}: PRMs trained
on human-written derivations transfer to model-written derivations only
when the model's CoT distribution is close to the labelled corpus.

\textbf{Why it matters for this week.} PRM800K is the canonical PRM
dataset and the empirical anchor for the PRM-vs-ORM debate in
section~6 of \texttt{notes.tex}. It also pre-figures the auto-labelling
race that produced Math-Shepherd and OmegaPRM by quietly noting that 800K
human step labels cost a small fortune.

\begin{quotation}
\itshape ``Process supervision provides a more precise reward signal than
outcome supervision and reduces the model's reliance on lucky end-state
matches. The benefit grows with the complexity of the reasoning chain
and the proportion of solutions in the training set that contain
partially correct work.''
\end{quotation}
The last clause is the crux: PRM advantage scales with the share of
``almost right'' chains. If your model is right or comprehensively wrong,
ORMs suffice. If most of its mistakes are last-step arithmetic on
otherwise correct derivations, a PRM is a much sharper teaching signal.

% ===========================================================================
\section*{Sourced from}

\begin{itemize}
  \item \texttt{research/llm\_rl\_sota.md} sections 3 (critic-free PG), 4 (RLVR), 5 (PRMs).
  \item \texttt{research/safety\_rl\_sota.md} section 5 (reward-hacking failure modes feeding into week 21).
\end{itemize}

\end{document}
